\documentclass[11pt,a4paper]{jsarticle}
\usepackage[utf8]{inputenc}
\usepackage{amsmath}
\usepackage{amssymb}
\usepackage{geometry}
\geometry{left=20mm,right=20mm,top=25mm,bottom=25mm}

\begin{document}

% --- タイトル部分 ---
\begin{center}
    {\LARGE \textbf{計測工学 第7回 実験1-光センサ}} \\
    \vspace{2mm}
    {\large 実験日：2026年5月28日}
\end{center}

\vspace{5mm}

% --- 1 目的と概要 ---
\section{目的と概要}
（※ここに目的と概要の本文を記述してください）

\vspace{3mm}

% --- 2 使用器具 ---
\section{使用器具}
班ごとに以下の器具を用意すること。また、実際に使用した器具の製造元、型番、シリアルナンバー等を記録すること。

\begin{itemize}
    \item 直流安定化電源
    \item ディジタルマルチメータ
    \item それらを繋ぐ線
    \item 簡易暗箱（LED付き）
    \item アルデュイーノ
\end{itemize}

簡易暗箱に付いているLEDライトはLDOが実装されており、自動車のバッテリーに接続して定格12Vで発光するように設計されている。実験時は安定化電源を用いて電源電圧 8--12V 前後の範囲で給電することで、箱内の照度を1000\,lx程度まで制御できる。

加えて、今回評価する光センサや必要な部材のセットを配布する。実験のために照度計を用意しているが、台数に限りがあるため複数の班で共有して使用すること。スマートフォンの照度センサを用いた照度計アプリがあるので、これらをインストールして用いても良い（レポートには班で使用した照度計、あるいはスマートフォンの型番やアプリの情報を記載すること）。また、実験室の室温を記録しておくこと。

\vspace{3mm}

% --- 3 実験課題 ---
\section{実験課題}
各課題で要求される実験データを取得できれば、具体的な方法（手順）はアレンジして構わない。レポートには実際に自分達が実施した手順や実験系をまとめること。実験課題1と2は必須、時間に余裕があるグループは課題3、4にも取り組むこと。

\subsection*{実験課題1. CdS の照度-抵抗値特性の評価}
照度を変化させた際のCdSの抵抗値の変化を評価し、グラフにまとめよ。また、抵抗値と照度の換算係数を求めよ。

\subsubsection*{参考手順（実験課題2と同時並行で実施せよ）}
\begin{enumerate}
    \item 暗箱のLEDを消灯している際の照度を測定する。
    \item CdS を暗箱中に入れ、上記照度の際の抵抗値を測定する。
    \item LEDを点灯し、照度を10\,lx程度に調整し、照度と CdS の抵抗値を測定する。
    \item 照度を10\,lx 刻みで100\,lxまで段階的に変化させ、その際の照度と CdSの抵抗値を測定する。
    \item 照度を100\,lx 刻みで1000\,lx まで段階的に変化させ、同様に照度とCdSの抵抗値を測定する。
\end{enumerate}

\subsection*{実験課題2. フォトトランジスタの照度-電流特性の評価}
照度を変化させた際のフォトトランジスタのコレクタ--エミッタ間電流（光電流$I_L$）を評価し、グラフにまとめよ。また線形領域のグラフの傾きから、光電流と照度の換算係数を求めよ。
\subsubsection*{参考手順（実際には実験課題1と同時並行で実施せよ）}
\begin{enumerate}
    \item CdS の評価回路に用いられているエミッタ抵抗の抵抗値を測定し、記録する。
    \item 実験課題1と同様に暗箱のLEDを消灯している際の照度を測定する。
    \item 測定回路を暗箱中に入れ、エミッタ抵抗の両端電圧を測定する（電源電圧はデータシート等をもとに適切に設定せよ）。
    \item LEDを点灯し、照度を10\,lx程度に調整し、同様に照度と電圧を測定する。
    \item 照度を10\,lx 刻みで100\,lx まで段階的に変化させ、同様に測定する。
    \item 照度を100\,lx刻みで1000\,lx まで段階的に変化させ、同様に測定する。
    \item 使用した抵抗の抵抗値と測定した電圧から電流を算出する。
\end{enumerate}

\subsection*{実験課題3. CdS を用いた照度測定}
実験室内の任意の箇所や、照明条件を変え、その時の照度と CdSの抵抗値を測定せよ。また、実験課題1の結果を用いて測定したCdSの抵抗値から照度を求め、照度計による測定値と比較せよ。
\subsection*{実験課題4. フォトトランジスタを用いた照度測定}
実験課題3と同様にフォトトランジスタの測定回路を用い、任意の環境におけるエミッタ抵抗の両端電圧と照度計による照度の測定値を記録せよ。また、実験課題2の結果を用い、測定した電圧から照度を求め、照度計による測定値と比較せよ。

\subsection*{追加課題}
実験課題1--4が全て終わり時間に余裕がある班は CdS とフォトトランジスタの応答性に関する以下の実験に取り組め。追加課題に関する実験結果および考察に取り組んだレポートには追加で加点を行う。

\subsubsection*{参考手順}
\begin{enumerate}
    \item オシロスコープを用意する。
    \item CdS およびフォトトランジスタによる照度測定回路を構成し、両回路の出力をオシロスコープに接続し観測する。
    \item 暗箱内に回路を入れ、LEDを点灯させた瞬間の両回路の立ち上がり特性を比較し、波形を記録する（フォトトランジスタの立ち上がりでシングルトリガを掛けると良い）。
    \item 同様にLEDを消灯させた際の両者の立ち下がり特性についても比較し、記録する。
\end{enumerate}

\vspace{3mm}

% --- 4 考察課題 ---
\section{考察課題}

\subsection{CdS の照度-抵抗値特性}
実験により得たグラフおよびをデータシートと比較し、実験結果について考察せよ。また、実験した各照度における CdS の抵抗値の実測結果と、実験で得た特性から計算される抵抗値の差をグラフにまとめ、考察せよ。

\subsection{フォトトランジスタの照度-電流特性}
実験により得たエミッタ抵抗の両端電圧とエミッタ抵抗の抵抗値から、フォトトランジスタの光電流を求め、照度-光電流のグラフを作成せよ。また、この結果を用い、照度と光電流が両対数軸上で線形であると仮定して光電流から照度を求める変換式を求めよ。

\subsection{不確かさの要因}
自分達が実施した各実験の手順や測定系について考察し、それぞれ測定値の不確かさの要因として考えられる項目を挙げよ。
\subsection{任意課題：不確かさの評価}
上記不確かさ要因に基づき、以下の項目について不確かさを評価し、バジェットシートをまとめよ。

\begin{enumerate}
    \item 実験課題1で求めた特性
    \item 実験課題2で求めた光電流と照度の換算式
    \item 実験課題3において、CdSの抵抗値と実験課題1で求めた特性から算出した実験室の照度
    \item 実験課題4において、フォトトランジスタ回路のエミッタ抵抗における電圧降下と実験課題2で求めた換算式から算出した実験室の照度
\end{enumerate}
※上記4つ全てに取り組む必要は無い（実施した実験内容にあわせて、考えやすい項目についてまとめること）。また、不確かさの要因は多数考えられるが、影響が大きそうで考えやすい項目に絞って構わない。

\vspace{3mm}

% --- 5 レポートについて ---
\section{レポートについて}
今回の実験について班で1通レポートを作成し、提出すること。本文の様式は自由であるが、レポートには以下の内容を記載すること。

\begin{description}
    \item[著者の情報] 特別な表紙は不要。1ページ目の上部に班員全員の氏名と学生番号を列挙すること。
    \item[目的] 実際に行なった実験や考察、結論と対応するように数行程度でまとめること。
    \item[実験方法] 自分達が行なった実験の手順や使用した器具について、実験課題ごとにまとめること。
    \item[実験結果] 各実験課題ごとに図表を用いて結果をまとめること。
    \item[考察] 各課題に対して考察をまとめること。
    \item[結論] 簡潔に結論をまとめること。
    \item[参考文献] レポートを書くにあたり参考にした文献があれば記載すること。
\end{description}

\end{document}